\section{Discussion and Engineering Implications}
\label{sec:discussion}

\subsection{What the Benchmark Says About Surrogate-Assisted Search}
\label{sec:boundary_dominance}
The benchmark separates three sources of gain over the pool best. The first is local exploitation: perturbing the 20 best pool designs raises the best range from \NvPool\ to \NvMeanLocTen\,km. The second is letting the search center move, which adds \NvDAdaptvLocTen\,km without any model. The third is the model. For the active-learning loop as originally implemented it adds nothing measurable over the adaptive local search ($\NvDsALvAdapt$\,km, interval $\NvCIALvAdapt$\,km); with the washout repair it adds \NvDALRvAdapt\,km; for TuRBO it adds \NvDTurbovAdapt\,km. A claim that a surrogate-assisted loop outperforms its initial sample therefore says little: here \NvShareAdapt\% of the gain of the best strategy is available without a surrogate.

Three practical points follow. The first concerns the initial sample. The pool of 9{,}800 designs took $97\%$ of the evaluations, yet TuRBO started from 1{,}000 designs ends \NpTurboLossB\ to \NpIndTurboVsFull\,km below its full-pool result, and from 288 of them it still exceeds the best pool design in every run (Section~\ref{sec:pool_size}). For a model-based search on this problem a pool a few times the search budget recovers all but about $\NpIndTurboVsFullPct\%$ of the range, and evaluations beyond that are better spent on the search. The other two concern the strategies themselves. A constraint that can be satisfied by construction should be repaired before evaluation in every strategy; leaving it to a classifier cost the published loop \NvWashAL\% of its budget and most of its advantage. And a local Gaussian-process model with a trust region was both simpler to run and more reliable than a global neural ensemble with a classifier: it produced no exhausted run, the highest feasible fraction, and the best mean. The global loop remains useful where a feasibility model is wanted for other purposes, and its result after repair is \NvPctTurbovALR\% below that of TuRBO, but the data do not favor it for this problem.

\subsection{Bounds, Constraint Fidelity and the Search Strategy}
\label{sec:decoupling_contributions}
Relaxing four bounds together raised the mean best range of a simple elite search from \NpaperBndMean\ to \NwideBndMean\,km (Section~\ref{sec:bound_sensitivity}), several times the spread between the three leading strategies in Table~\ref{tab:budget}. The single-bound runs attribute this mainly to battery mass: $0.2$\,kg more battery gives \NbpBmassGain\,km, more than the difference between TuRBO and the adaptive local search. Suo et al.\ report the analogous result for a fuel-powered BWB UAV, where the fuel mass dominates the sensitivity of range~\cite{suo2026kriging}. The best designs sit at the battery-mass limit and close to the lowest permitted cruise speed, so the result of any strategy is set largely by those two requirements, and none of the designs should be described as a global optimum of the vehicle.

The fidelity of the stability constraint matters in a different way. The Euler calibration does not change the ranking of the strategies, but it moves the effective upper limit on the lattice margin from $0.25$ to \NcalCapEquiv, and the best design of \NcalOutLocTen\ of \NcalOutNLocTen\ static-local runs, \NcalOutAlR\ of \NcalOutNAlR\ AL-R runs, \NcalOutAdapt\ of \NcalOutNAdapt\ adaptive-local runs and \NcalOutTurbo\ of \NcalOutNTurbo\ TuRBO runs falls outside the calibrated band. A search that pushes against a constraint evaluated by a low-order model will find the model's error, and the calibrated limit should be used in the constraint from the start in future work.

The chord clamp of the loft is a validity guard. Inside the working bounds it is never active, and with widened bounds a search with the clamp and one without it reach the same range within the uncertainty of ten paired seeds. The loft does not make the center-body chord monotone (Section~\ref{sec:loft_def}); the excursions are below six millimeters and have no aerodynamic significance at this fidelity, but the earlier description of the loft as monotone by construction was incorrect.

\subsection{Limitations}
\label{sec:limitations}
\begin{itemize}
    \item \textbf{One aircraft, one pipeline, one budget.} All conclusions about search strategies hold for this 14-variable problem, this evaluator and 287 further evaluations. The ordering was checked for pools of 288 to 9{,}800 designs drawn from one Latin hypercube sample, for ten independent pools of 1{,}000 designs (Section~\ref{sec:pool_size}), and for budgets of 50 to 287 evaluations (Table~\ref{tab:anytime}); on the independent pools the lead of TuRBO over the adaptive local search has the same sign but does not reach the 5\% level. At twice the budget the ordering is unchanged and the gaps are smaller (Table~\ref{tab:long}). A more expensive evaluator, which would allow fewer evaluations, could change the ordering, and a recent benchmark of aerodynamic shape optimizers over 103 tasks found that optimizer rankings vary substantially between problem classes~\cite{fazliani2026shapebench}, so the ordering found here should not be assumed for other aircraft problems.
    \item \textbf{No viscous three-dimensional validation and no experiment.} The lattice assumes attached inviscid flow; at the cruise condition of the best designs ($C_L \approx 0.25$, $\alpha \approx 4^\circ$ in the lattice) this is plausible but unverified. Profile drag is a strip estimate with an assumed transition location, and AeroBuildup, whose section drag comes from NeuralFoil, gives an $L/D$ higher by a median factor of $1.46$. The Euler solutions check lift and the static margin only, on meshes that are not converged to better than about $5\%$ in lift, and contain no boundary layer. Absolute range and the true static margin require RANS solutions or wind-tunnel and flight data.
    \item \textbf{Stability constraint.} The optimization used the lattice static margin, which the Euler solutions show to be low by about $0.035$; the correction was confirmed on \NoosN\ designs outside the fit but only against Euler solutions. The calibrated constraint was applied after the fact and not inside a search. Lateral-directional stability, control power, aeroelastic effects and structural sizing are not modeled.
    \item \textbf{Mass and energy model.} The airframe mass is a linear function of reference area, the propulsive efficiency is an assumed four-point map, and the usable energy assumes $80\%$ depth of discharge at a steady $0.28\,C$ discharge. Range is linear in the usable energy, so these assumptions rescale all designs equally and do not affect the comparisons.
    \item \textbf{Fixed battery centroid and bay packaging.} The battery centroid is held at $0.82\,c_{r}$ for every design; it is never optimized and is not checked against the bay geometry. The packaging check integrates a center-body volume along the span with $45\%$ of the local chord as usable width and does not constrain where the battery lies along the chord, so a detailed internal layout may not realize the assumed position. Section~\ref{sec:robustness} shows that the best designs tolerate a shift of $\pm 0.04\,c_r$.
    \item \textbf{Guards inside the evaluator.} The best designs lie at the battery-mass bound and close to the evaluator's span-efficiency guard ($e \le 1$), which is a numerical safeguard and not a physical limit for a wing with winglets. The optima reported here are optima of this evaluator including its guards.
    \item \textbf{Benchmark implementation.} The active-learning loop and the baselines are single implementations with untuned hyperparameters. Trust-region methods that model the constraints themselves, such as SCBO~\cite{eriksson2021scbo} and FuRBO~\cite{ascia2025furbo}, were not tested; they need constraint values, and the strategies compared here see only whether a design is feasible. Two implementation issues that affected the comparison were found and corrected during this study (Section~\ref{sec:budget_protocol}); others may remain. The code and all run histories are released so that the comparison can be checked.
    \item \textbf{Statistics.} Ten to twenty runs per strategy detect differences of several kilometers; smaller differences between the leading strategies cannot be resolved, and the confidence intervals state what remains possible.
\end{itemize}
